\subsection{Reinforcement-Learning Allocation Policy}\label{subsec:rl}
\textit{(Aldo Patrone)}\quad
The allocation policy is learned as a deep-RL agent, faithful to Middelhuis et
al.~\cite{middelhuis2025rl}. The environment is a discrete-event simulation of the process, the ``simulator'' the paper trains in, calibrated to BPIC-17 through our own engines
(interarrival gaps from the arrival engine, per-activity processing and waiting times and
BPMN-valid case sequences from the process model and branching engine). The Markov decision
process matches the paper (Appendix Fig.~\ref{fig:rl_mdp}):

\begin{itemize}
  \item \textbf{State} ($2|R|+|A|$): for every resource a busy flag $\delta_i\in\{0,1\}$ and,
  when busy, which activity it executes $\eta_i\in[0,1]$. For every activity the normalized
  queue of unassigned instances $\kappa_j=\min(|K_{a_j}|/100,\,1)$.
  \item \textbf{Action} ($|R|\cdot|A|+1$): every (resource, activity) assignment plus a
  \emph{postpone}. Infeasible pairs are removed by an \emph{action mask}: a pair is infeasible
  when there is no waiting instance of that activity, or the resource is busy, unavailable or
  unauthorized.
  \item \textbf{Reward}: the objective is to minimize the mean cycle time. Each time the
  simulation advances by $\Delta t$ the agent receives $-\Delta t\cdot n_{\text{active}}/3600$,
  which telescopes to the negative total cycle time (Def.~13--14 of~\cite{middelhuis2025rl}),
  plus a small potential-based shaping on the load Gini for fairness.
\end{itemize}
We train the paper's algorithm, PPO~\cite{schulman2017ppo}, in its action-masking variant
\emph{MaskablePPO}~\cite{huang2022masking} (sb3-contrib~\cite{raffin2021sb3}), which handles the
variable feasible-action set through the mask. Postponing accrues the same time-in-system penalty, so
idling is discouraged without an artificial cost. A bounded fallback forces the least-loaded
assignment after repeated postpones to keep training stable.

The two learners optimize the standard objectives. MaskablePPO maximizes the clipped PPO surrogate,
with the action mask zeroing infeasible actions before normalization, and REINFORCE follows the
policy-gradient rule~\cite{williams1992reinforce, sutton2018rl} with a return-to-go advantage and a
running baseline.

The learned policy is deployed onto the integrated simulator through the swappable
\emph{AllocationStrategy.pick(candidates, context)} interface. Because the engine allocates one
ready activity instance at a time, deployment is a \emph{per-task reduction}: the mask is
restricted to the current activity, so the policy selects the resource for the task at hand (a
postpone returns \emph{None} and the core suspends and retries). The state's $\delta$ and
$\eta$ are reconstructed from the live \emph{ResourceEngine}. The global queue term $\kappa$ is
not observable at a single decision, so it is reduced to the locally known waiting instance,
a disclosed reduced-observation deployment, while the full $\kappa$ evaluation is carried out in
the paper-style environment itself.

As a lightweight counterpart to MaskablePPO we implemented, from scratch in pure \emph{numpy}, a
per-candidate-scoring REINFORCE~\cite{williams1992reinforce, sutton2018rl} policy. For a
like-for-like comparison it is trained in the \emph{same} paper-faithful environment as PPO
(identical dynamics and reward), acting through the per-task decision interface, so both learned
policies are scored on the same footing in Sec.~\ref{ssub:rl-eval}. A
fairness-weight sensitivity over $w_{\mathrm{fair}}\in\{0,1,5\}$ (Appendix
Fig.~\ref{fig:alloc_cmp_plot}) tests whether any fairness difference is a real algorithmic property
rather than an artifact of the reward weight. The numpy policy serializes to JSON so its deployment
needs no machine-learning dependency at runtime, and its hyperparameters are in
Table~\ref{tab:rl_hparams}.
